
% =====================================================================
\section{Finding: clustering and next-day liquidity}
\label{sec:liquidity}
% =====================================================================

This is the question the paper leaves open. Clustering measures noise-trader
activity; what does that activity do to the cost of trading?

\subsection{Specification, and why it is the dynamic one}

\begin{table}[htbp]
\centering
\caption{Price clustering and next-day liquidity}
\label{tab:rq2}
\small
\begin{threeparttable}
\begin{tabular}{@{}lrrrrr@{}}
\toprule
Outcome & Predictive & SE & Dynamic & SE & Stock-days \\
\midrule
Effective spread (log) & 0.0367*** & (0.0080) & 0.0188** & (0.0075) & 74,034 \\
Price impact, 60s (bp) & 0.1587 & (0.1476) & 0.0417 & (0.1498) & 74,034 \\
Realized spread, 60s (bp)$^{\dagger}$ & -0.0039 & (0.1705) & 0.0776 & (0.1691) & 74,034 \\
Depth at best quote (log)$^{\dagger}$ & -0.0616*** & (0.0092) & -0.0441*** & (0.0085) & 74,034 \\
Realised variance (log)$^{\dagger}$ & -0.0114 & (0.0187) & -0.0065 & (0.0177) & 74,030 \\
Variance-ratio deviation$^{\dagger}$ & 0.0106** & (0.0050) & 0.0106** & (0.0050) & 74,032 \\
Amihud illiquidity (log)$^{\dagger}$ & 0.0372 & (0.0301) & 0.0305 & (0.0317) & 69,491 \\
\bottomrule
\end{tabular}
\begin{tablenotes}[flushleft]\footnotesize
\item Each cell is the coefficient on $M^{BLarge0}_{t-1}$, the round-price share of buy-initiated large-trade volume. All specifications carry stock and day fixed effects and the full control set (relative tick size, lagged log turnover, lagged log realised variance, lagged log effective spread, the overnight return, and the opening-digit by low-volatility interactions, timed to match the regressor: lagged interactions beside $M_{t-1}$, same-day ones beside $M_t$). The dynamic column adds the outcome's own lag, so its coefficient is the incremental predictive content of clustering given the outcome's own one-day lag; Section~\ref{sec:robust} probes deeper lags. Standard errors are clustered by stock and by day. Days on the 0.1-yen grid are excluded, following Ohta (2026). $^{\dagger}$ secondary outcome, not pre-specified. $^{*}p<0.1$, $^{**}p<0.05$, $^{***}p<0.01$. These are predictive associations in a single quarter, not causal estimates.
\end{tablenotes}
\end{threeparttable}
\end{table}


Three specifications are estimated for each outcome and
Table~\ref{tab:rq2} reports two of them. All carry stock and day fixed
effects, the full control set, and standard errors clustered on both margins.

The \emph{predictive} column regresses today's liquidity on yesterday's
clustering. It fixes the ordering but not much else: both series are persistent,
so a stock that is illiquid and clustered today was probably both yesterday, and
the coefficient partly recovers that.

The \emph{dynamic} column adds the outcome's own lag, and it is the one the
argument rests on. Its coefficient answers a sharper question --- does
yesterday's clustering say anything about today's liquidity that the
outcome's \emph{own one-day lag} did not already say? With 78 time
periods the resulting dynamic-panel bias is of order $1/T$
\parencite{Nickell1981} and can be neglected. One lag is not the same as
``the outcome's history'', so Section~\ref{sec:robust} re-asks the question
against five lags.

\subsection{What the panel says}

The coefficient on lagged clustering in the effective-spread regression is
0.0188, significant at the 5\% level. Read economically, a one-standard-deviation
increase in the round-price share of large-trade volume is associated with an
effective spread about 0.2\% higher the following day. For the price-impact outcome the
coefficient is
0.0417, not statistically distinguishable from zero. Depth at the best quote moves the
complementary way (-0.0441, significant at the 1\% level): the day after heavy
round-price trading quotes are wider and thinner.

\paragraph{Which way the open question falls.}
Section~\ref{sec:intro} set out two incompatible predictions. On the
\textcite{Kyle1985} reading, more noise trading means a deeper market and
cheaper trading. On the stale-order reading, clustering marks liquidity supplied
by participants about to be picked off, and trading should get more expensive.
The sign here is positive, favouring the second reading: days following heavy round-price trading are days when trading costs more. That is one sample of one market, and it is
reported as an association rather than a mechanism --- and where the coefficient
is not distinguishable from zero, as a direction the data do not resolve.

\subsection{Which measure carries the information}

\begin{table}[htbp]
\centering
\caption{Which clustering measure carries the information?}
\label{tab:rq2alt}
\small
\begin{threeparttable}
\begin{tabular}{@{}lrrr@{}}
\toprule
Regressor (lagged) & Coefficient & SE & Stock-days \\
\midrule
$M^{BLarge0}$ & 0.0192*** & (0.0075) & 74,034 \\
$M^{SLarge0}$ & 0.0249*** & (0.0086) & 74,090 \\
$M^{BSmall0}$ & 0.0213** & (0.0093) & 74,218 \\
$M^{SSmall0}$ & 0.0154 & (0.0108) & 74,218 \\
$M^{0}$ & 0.0416*** & (0.0123) & 74,218 \\
\bottomrule
\end{tabular}
\begin{tablenotes}[flushleft]\footnotesize
\item Outcome is the log effective spread; the dynamic specification of Table~\ref{tab:rq2} with the clustering measure replaced. The mechanism concerns large trades executing against stale round-price limit orders, so the large-trade measures are where an effect should appear.
\end{tablenotes}
\end{threeparttable}
\end{table}


Table~\ref{tab:rq2alt} substitutes the other clustering measures. If the
relationship works through the channel Ohta describes --- large orders taking
stale round-price limit orders --- it should live in the large-trade measures and
not in the small-trade ones. This is a weaker test than
the placebo in Section~\ref{sec:robust}, but it uses the paper's own internal
contrast, and it deserves an explicit reading rather than a table alone.

The contrast the mechanism predicts does not appear. Per standard deviation of the sorting measure, the small-trade share moves next-day spreads by 0.16\% against the large-trade share's 0.23\% --- comparable, not smaller --- and the pooled measure $M^{0}$ carries the largest coefficient of all (0.0418, $t=3.41$). Elevated round-price trading of any kind predicts wider spreads; the specific stale-limit-order channel, which lives in large trades, is not what separates the measures here. That reads as a caution against the mechanism interpretation of the daily association, not against the association itself.

\paragraph{Verdict.}
Clustering carries a measurable association with next-day liquidity beyond the outcome's own one-day lag, in a specification that absorbs every stock-level and market-wide effect. Section~\ref{sec:robust} bounds how hard this can be leaned on: the pooled estimate is significant at the 5\% level, but it falls below conventional significance once the outcome's history means five lags rather than one; its sign reverses between the within-stock and cross-sectional designs. The association is descriptive: no experiment, and no claim that intervening on clustering would move spreads.
